\PassOptionsToPackage{table,xcdraw}{xcolor}
\documentclass[10pt, a4paper]{article}

\usepackage[T1]{fontenc}
\usepackage[utf8]{inputenc}
\usepackage{tgpagella}
\usepackage{tgcursor}
\usepackage[spanish, es-noshorthands]{babel}
\usepackage{amsmath, amssymb}
\usepackage{booktabs}
\usepackage{tabularx}
\usepackage{xcolor}
\usepackage{tikz}
\usepackage{pgfplots}
\pgfplotsset{compat=1.18}
\usepackage[most]{tcolorbox}
\usepackage[margin=1.6cm]{geometry}
\usepackage{fancyhdr}
\usepackage{lastpage}

\setlength{\headheight}{16pt}
\setlength{\emergencystretch}{2em}

\definecolor{ucbblue}{RGB}{0, 51, 102}
\definecolor{siggreen}{RGB}{0, 105, 55}
\definecolor{neutralbg}{RGB}{246, 248, 251}
\definecolor{greenbg}{RGB}{242, 250, 245}

\pagestyle{fancy}
\fancyhf{}
\lhead{\textbf{IMT-221 --- Circuitos Electrónicos II}}
\chead{Solucionario Autónomo --- Práctica P (Divulgable)}
\rhead{W17-S03 -- 27/11/2026}
\cfoot{\footnotesize Página \thepage\ de \pageref{LastPage} \quad|\quad Artefacto A2 --- Solucionario Práctica P (NO contiene Examen E)}
\renewcommand{\headrulewidth}{0.5pt}
\renewcommand{\footrulewidth}{0.4pt}

\begin{document}

\begin{center}
    {\large\textbf{\textcolor{ucbblue}{Universidad Católica Boliviana ``San Pablo'' --- Sede Tarija}}}\\[0.08cm]
    {\Large\textbf{Solucionario Autónomo --- Práctica P (Divulgable a Estudiantes)}}\\
    {\normalsize IMT-221: Circuitos Electrónicos II $\cdot$ Preparación para Recuperatorio Extraordinario EC2}
\end{center}

\begin{tcolorbox}[colback=greenbg, colframe=siggreen, boxrule=0.7pt, arc=1.5mm, left=6pt, right=6pt, top=4pt, bottom=4pt]
    \footnotesize
    \textbf{Uso estudiantil:} Consultar después de resolver la Práctica P (A1). \textbf{Este documento corresponde exclusivamente al Caso P y no contiene enunciados ni respuestas del Examen E.} Todos los valores y gráficas son analíticos/ilustrativos, no simulaciones ejecutadas ni mediciones físicas.
\end{tcolorbox}

% =========================================================================
% SOLUCIÓN P1 (25 PTS)
% =========================================================================
\section*{P1 --- Sensor diferencial e In-Amp (25 puntos)}

\begin{enumerate}
    \setlength{\itemsep}{2pt}
    \item $v_d(t) = v_+(t) - v_-(t) = \bigl(0.006 - (-0.006)\bigr)\sin(2\pi\cdot 10\,t) = 0.012\sin(2\pi\cdot 10\,t)\,\mathrm{V}$ ($12\,\mathrm{mV_p}$); \quad $v_{cm}(t) = \dfrac{0.800 + 0.800}{2} = 0.800\,\mathrm{V}$.
    \item $v_o(t) = A_d v_d(t) + A_{cm} v_{cm}(t) = 100\bigl(0.012\sin(2\pi\cdot 10\,t)\bigr) + 0.010(0.800) = 1.20\sin(2\pi\cdot 10\,t) + 0.008\,\mathrm{V}$.
    \item $\mathrm{CMRR} = \dfrac{100}{0.010} = 10\,000$; \qquad $\mathrm{CMRR}_{\mathrm{dB}} = 20\log_{10}(10\,000) = 80\,\mathrm{dB}$.
    \item La salida contiene un \textit{offset} común de $+8\,\mathrm{mV}$; no se puede tratar como una senoide centrada en $0\,\mathrm{V}$ sin declarar compensación o separación de DC antes de analizar la etapa siguiente.
\end{enumerate}

% =========================================================================
% SOLUCIÓN P2 (25 PTS)
% =========================================================================
\section*{P2 --- Etapa no inversora y limitaciones dinámicas (25 puntos)}

\begin{enumerate}
    \setlength{\itemsep}{2pt}
    \item $A_v = 1 + \dfrac{R_f}{R_g} = 1 + \dfrac{10\,\mathrm{k\Omega}}{10\,\mathrm{k\Omega}} = 2$; \quad $V_{o,p} = 2(1.20\,\mathrm{V_p}) = 2.40\,\mathrm{V_p} < 3.0\,\mathrm{V}$ (cumple excursión lineal $\pm 3.0\,\mathrm{V}$).
    \item $\mathrm{NG} = 2 \implies f_{BW} \approx \dfrac{f_T}{\mathrm{NG}} = \dfrac{100\,\mathrm{kHz}}{2} = 50\,\mathrm{kHz}$. Tanto $10\,\mathrm{kHz}$ como $20\,\mathrm{kHz}$ están por debajo de $50\,\mathrm{kHz}$.
    \item $\mathrm{SR}_{\mathrm{req}}(10\,\mathrm{kHz}) = 2\pi(10\,000)(2.40) = 150\,796\,\mathrm{V/s} \approx 0.151\,\mathrm{V/\mu s}$ (cumple $\le 0.20\,\mathrm{V/\mu s}$); \quad $\mathrm{SR}_{\mathrm{req}}(20\,\mathrm{kHz}) = 2\pi(20\,000)(2.40) = 301\,593\,\mathrm{V/s} \approx 0.302\,\mathrm{V/\mu s}$ (incumple $> 0.20\,\mathrm{V/\mu s}$). La operación nominal a $10\,\mathrm{Hz}$ requiere $\mathrm{SR}_{\mathrm{req}}(10\,\mathrm{Hz}) \approx 0.000151\,\mathrm{V/\mu s}$.
    \item Solo \textbf{$10\,\mathrm{kHz}$ pasa} simultáneamente los tres chequeos ($20\,\mathrm{kHz}$ falla por SR). Cumplir $f < f_{BW}$ (pequeña señal) no garantiza ausencia de distorsión en gran señal ni desempeño exacto del op-amp real.
\end{enumerate}

% =========================================================================
% SOLUCIÓN P3 (25 PTS)
% =========================================================================
\section*{P3 --- Rectificación convencional y de precisión (25 puntos)}

\begin{enumerate}
    \setlength{\itemsep}{2pt}
    \item Para $0.20\,\mathrm{V_p}$: convencional $\max(0.20 - 0.50, 0) = 0\,\mathrm{V}$; precisión $0.20\,\mathrm{V_p}$. Para $0.80\,\mathrm{V_p}$: convencional $0.30\,\mathrm{V_p}$; precisión $0.80\,\mathrm{V_p}$.
    \item En semiciclo positivo ($v_i > 0$): el diodo del superdiodo conduce, la realimentación se cierra, $v_{o,P} \approx v_i$ y $v_a \approx v_{o,P} + v_D$. En semiciclo negativo ($v_i < 0$): el diodo bloquea, el lazo se abre, $v_{o,P} \approx 0$ por $R_L$ y $v_a$ tiende hacia el riel negativo.
    \item Referencia ideal $v_{\mathrm{ref}} = [0,\; 0,\; 0.2,\; 0.6,\; 0.8]\,\mathrm{V}$. Convencional: suma de errores $0 + 0 + 0.2 + 0.5 + 0.5 = 1.2\,\mathrm{V} \implies E_{\mathrm{MAE},D} = 1.2/5 = 0.240\,\mathrm{V}$. Precisión: suma $0 + 0 + 0.01 + 0.02 + 0.02 = 0.05\,\mathrm{V} \implies E_{\mathrm{MAE},P} = 0.010\,\mathrm{V}$ (errores discretos ilustrativos, no ensayos reales).
    \item La tensión interna negativa $v_a$ mantiene bloqueado al diodo $D_2$ con el lazo abierto; no produce una salida rectificada negativa porque $D_2$ impide corriente hacia $R_L$.
\end{enumerate}

\begin{tcolorbox}[colback=neutralbg, colframe=ucbblue!60, boxrule=0.5pt, arc=1mm, left=6pt, right=6pt, top=3pt, bottom=3pt]
    \scriptsize\textbf{Algoritmo descriptivo reproducible de $E_{\mathrm{MAE}}$ (solo verificación analítica, sin LTspice):}\\
    \texttt{para cada muestra k: vref[k] = max(vi[k], 0); errD[k] = abs(voD[k] - vref[k]); errP[k] = abs(voP[k] - vref[k])}\\
    \texttt{EMAE\_D = suma(errD)/N; \quad EMAE\_P = suma(errP)/N}
\end{tcolorbox}

% =========================================================================
% SOLUCIÓN P4 (25 PTS) + V6 BODE ANALÍTICO + V8 AUTOEVALUACIÓN
% =========================================================================
\newpage
\section*{P4 --- Filtro pasa-bajos Sallen--Key de segundo orden (25 puntos)}

\begin{enumerate}
    \setlength{\itemsep}{2pt}
    \item $RC = 16\,000 \times 10^{-6} = 0.016\,\mathrm{s}$; $H(s) = \dfrac{1}{1 + 0.032s + 0.000256s^2} = \dfrac{1}{(1 + 0.016s)^2}$. Entonces $\omega_0 = \dfrac{1}{0.016} = 62.5\,\mathrm{rad/s}$, $f_0 = \dfrac{62.5}{2\pi} \approx 9.947\,\mathrm{Hz}$ y $Q = 0.5$.
    \item $f_{-3\mathrm{dB}} = f_0\sqrt{\sqrt{2} - 1} \approx 9.947 \times 0.643594 \approx 6.402\,\mathrm{Hz} < f_0$ (en $Q = 0.5$, el corte a $-3\,\mathrm{dB}$ ocurre por debajo de $f_0$, a diferencia de Butterworth donde $Q = 1/\sqrt{2}$).
    \item En $f = f_0$: $|H(j\omega_0)| = 0.5$, $G_{\mathrm{dB}}(f_0) = 20\log_{10}(0.5) \approx -6.0206\,\mathrm{dB}$, $\varphi(f_0) = -90^\circ$. Con la senoide de prueba de $0.50\,\mathrm{V_p}$: $V_{o,p}(f_0) = 0.50 \times 0.50 = 0.25\,\mathrm{V_p}$.
    \item A $50\,\mathrm{Hz}$: $x = 50 / 9.947 \approx 5.02655$, $|H(50\,\mathrm{Hz})| = \dfrac{1}{1 + x^2} \approx 0.03807$, $G_{\mathrm{dB}}(50\,\mathrm{Hz}) \approx -28.39\,\mathrm{dB}$. \textbf{No} satisface $-40\,\mathrm{dB}$ a $50\,\mathrm{Hz}$ (ni constituye evidencia física).
\end{enumerate}

% V6 — Bode analítico calculado para el filtro P (R=16 kOhm, C=1 uF, f0=9.947 Hz)
\begin{center}
\noindent
\begin{minipage}[t]{0.48\textwidth}
\centering
\begin{tikzpicture}
    \begin{axis}[
        width=0.98\textwidth, height=4.4cm,
        xmode=log, xmin=1, xmax=100, ymin=-42, ymax=2,
        xlabel={\scriptsize Frecuencia $f$ [Hz]}, ylabel={\scriptsize Magnitud $G_{\mathrm{dB}}$ [dB]},
        grid=major, tick label style={font=\scriptsize},
        title={\scriptsize\textbf{V6a: Magnitud (calculada analíticamente)}}
    ]
    \addplot[ucbblue, very thick, domain=1:100, samples=120] {-20*log10(1 + (x/9.947)^2)};
    \addplot[only marks, mark=*, mark size=1.8pt, red!75!black] coordinates {(6.402, -3.01) (9.947, -6.02) (50, -28.39)};
    \node[anchor=south west, font=\tiny] at (axis cs:6.4, -3.0) {$f_{-3\mathrm{dB}}\!\approx\!6.40\,\mathrm{Hz}$};
    \node[anchor=north east, font=\tiny] at (axis cs:9.5, -7.2) {$f_0\!\approx\!9.95\,\mathrm{Hz}$ ($-6.02\,\mathrm{dB}$)};
    \node[anchor=east, font=\tiny] at (axis cs:46, -28.4) {$50\,\mathrm{Hz}$ ($-28.39\,\mathrm{dB}$)};
    \end{axis}
\end{tikzpicture}
\end{minipage}
\hfill
\begin{minipage}[t]{0.48\textwidth}
\centering
\begin{tikzpicture}
    \begin{axis}[
        width=0.98\textwidth, height=4.4cm,
        xmode=log, xmin=1, xmax=100, ymin=-185, ymax=5,
        xlabel={\scriptsize Frecuencia $f$ [Hz]}, ylabel={\scriptsize Fase $\varphi$ [$^\circ$]},
        ytick={0, -45, -90, -135, -180},
        grid=major, tick label style={font=\scriptsize},
        title={\scriptsize\textbf{V6b: Fase (calculada analíticamente)}}
    ]
    \addplot[siggreen, very thick, domain=1:100, samples=120] {-2*atan(x/9.947)};
    \addplot[only marks, mark=*, mark size=1.8pt, red!75!black] coordinates {(9.947, -90)};
    \draw[dashed, gray] (axis cs:9.947, -185) -- (axis cs:9.947, -90) -- (axis cs:1, -90);
    \node[anchor=south west, font=\tiny] at (axis cs:10.5, -88) {$f_0\!\approx\!9.95\,\mathrm{Hz}$ ($-90^\circ$)};
    \end{axis}
\end{tikzpicture}
\end{minipage}
\par\vspace{2pt}{\scriptsize\textit{Curvas calculadas analíticamente a partir de $H(s) = 1/(1+0.016s)^2$ ($Q=0.5$); no son simulaciones LTspice ni mediciones.}}
\end{center}

% V8 — Autoevaluación y errores frecuentes
\begin{tcolorbox}[colback=neutralbg, colframe=ucbblue, title=\textbf{V8: Autoevaluación y Errores Frecuentes de la Práctica (sin nota sumativa)}]
    \footnotesize
    \begin{itemize}\setlength\itemsep{2pt}
        \item[$\square$] ¿Se confundió $v_{cm}$ con $v_d$? Verificar resta ($v_+ - v_-$) y semisuma ($(v_+ + v_-)/2$). \textit{[Modelo / álgebra]}
        \item[$\square$] ¿Se comprobó $\mathrm{SR}$ en $\mathrm{V/\mu s}$ y excursión independientemente de $f_{BW}$? \textit{[Unidades / modelo]}
        \item[$\square$] ¿Se usó el cortocircuito virtual con el diodo bloqueado? Eso es incorrecto. \textit{[Modelo de lazo]}
        \item[$\square$] ¿Se asumió que $-3\,\mathrm{dB}$ ocurre en $f_0$ con $Q = 0.5$? Eso es incorrecto ($-6.02\,\mathrm{dB}$). \textit{[Interpretación]}
        \item[$\square$] ¿Se interpretó la prueba sinusoidal de Bode como predicción directa del pico de la media onda rectificada? Eso no es correcto. \textit{[Interpretación]}
    \end{itemize}
\end{tcolorbox}

\end{document}
